% GENERATED FILE. Edit canonical lessons, then run npm run generate:books.
% canonical-source-hash: fnv1a64:bb23afa959d2414f
% canonical-lessons: SA-C57-shaknoti, SA-C57-icchati

\chapter{Able, and Willing}
\label{ch:sa-able-and-willing}
\hlchaptermodality{57}

\begin{chapteropening}
\textbf{By the end of this chapter:} \emph{I can say what somebody is able to do and what they want to do, by putting the -\sk{तुम्} form in front of \sk{शक्नोति} or \sk{इच्छति}.}

Complete SA-C57-icchati and say both frames — the boy can read, and the boy wants to read.
\end{chapteropening}

\section[śaknoti]{\sk{शक्नोति} (śaknoti) — is able}

\label{lesson:SA-C57-shaknoti}



\begin{quote}
Build the past of \textbf{\sk{पठति}}. Then say \textbf{\sk{गन्तुम्}} and what its ending does.
\end{quote}

\subsection*{You'll want to know: \sk{शक्नोति}}
\textbf{\sk{शक्नोति}} (\emph{śaknoti}) — \textquotedblleft{}is able, can\textquotedblright{}.

Until now this book could say what somebody does and, since last chapter, what they did. It could not say what they \textbf{could} do.

\textbf{\sk{शक्नोति}} takes the \textbf{-\sk{तुम्}} form as its object — which is why it arrives two chapters after that ending and not before:

\begin{quote}
\textbf{\sk{बालः} \sk{गन्तुं} \sk{शक्नोति}}
\end{quote}

\begin{quote}
\emph{bālaḥ gantuṁ śaknoti}
\end{quote}

\begin{quote}
\textbf{the boy can go}
\end{quote}

Read the Sanskrit literally and it says \emph{the boy is-able to-go}. English puts \emph{can} in front of a bare verb; Sanskrit puts \textbf{\sk{शक्नोति}} at the end after an infinitive. Same idea, opposite order, and the infinitive is the piece that makes it possible.

You already know a word from this root without knowing it: \textbf{\sk{शक्ति}} (\emph{śakti}), \textquotedblleft{}power, capacity\textquotedblright{} — the same \textbf{\sk{शक्}}, made into a noun. English borrowed that one whole.

(The \textbf{\sk{म्}} of \textbf{\sk{गन्तुम्}} is written \textbf{\sk{ं}} here — \emph{gantuṁ} — because a following consonant softens it. You have read that mark since the asking verb.)

\begin{grammarlens}[title={What you've built}]
Ability: what somebody is able to do, not merely what they do.
\end{grammarlens}

\subsection*{Your turn}
\begin{itemize}
\raggedright
  \item \emph{Say it:} \emph{śaknoti}
  \item \emph{Say it:} \emph{bālaḥ gantuṁ śaknoti}
  \item \emph{Say it:} the same sentence with \emph{paṭhitum} instead, and say what it now means
  \item \emph{Read it:} \textbf{\sk{शक्नोति}}, and name the noun English borrowed from its root
\end{itemize}

\begin{tcolorbox}[breakable,colback=teal!4,colframe=teal!35!black,title={Before you move on}]
What form does \textbf{\sk{शक्नोति}} take as its object? (The \textbf{-\sk{तुम्}} form.) Say \textquotedblleft{}the boy can go\textquotedblright{}.
\end{tcolorbox}
\section[icchati]{\sk{इच्छति} (icchati) — wants}

\label{lesson:SA-C57-icchati}



\begin{quote}
Say \textquotedblleft{}the boy can go\textquotedblright{}. Which of the two words in it is the \textbf{-\sk{तुम्}} form?
\end{quote}

\subsection*{You'll want to know: \sk{इच्छति}}
\textbf{\sk{इच्छति}} (\emph{icchati}) — \textquotedblleft{}wants, wishes\textquotedblright{}.

There is nothing new to learn about the frame. It is the frame you built one lesson ago, with a different verb at the end of it:

\begin{quote}
\textbf{\sk{बालः} \sk{पठितुम्} \sk{इच्छति}}
\end{quote}

\begin{quote}
\emph{bālaḥ paṭhitum icchati}
\end{quote}

\begin{quote}
\textbf{the boy wants to read}
\end{quote}

Compare the pair and the shape holds still while the meaning turns:

\noindent
\begin{tabularx}{\linewidth}{@{}>{\raggedright\arraybackslash}X>{\raggedright\arraybackslash}X@{}}
\toprule
\textbf{sentence} & \textbf{meaning} \\
\midrule
\textbf{\sk{बालः} \sk{पठितुम्} \sk{शक्नोति}} & the boy \textbf{can} read \\
\textbf{\sk{बालः} \sk{पठितुम्} \sk{इच्छति}} & the boy \textbf{wants to} read \\
\bottomrule
\end{tabularx}

That is what makes the \textbf{-\sk{तुम्}} form worth the lesson it got two chapters back. It is not one construction; it is a socket, and ability and desire are the first two things to plug into it.

The root is \textbf{\sk{इष्}}, \textquotedblleft{}to seek\textquotedblright{} — and the noun it makes is one you may have met in English writing about India: \textbf{\sk{इच्छा}} (\emph{icchā}), \textquotedblleft{}wish, desire\textquotedblright{}.

\begin{grammarlens}[title={What you've built}]
Desire, in the frame you already had.
\end{grammarlens}

\subsection*{Your turn}
\begin{itemize}
\raggedright
  \item \emph{Say it:} \emph{bālaḥ paṭhitum icchati}
  \item \emph{Contrast:} the same sentence with \emph{śaknoti}, and say what changed
  \item \emph{Read it:} \textbf{\sk{इच्छति}}, then say what it means
\end{itemize}

\begin{tcolorbox}[breakable,colback=teal!4,colframe=teal!35!black,title={Before you move on}]
Say \textquotedblleft{}the boy wants to read\textquotedblright{}, then say \textquotedblleft{}the boy can read\textquotedblright{}.
\end{tcolorbox}
